\chapter{Stage 3b: The Retrieval-Utility Classifier}
\label{ch:ruc}

This is the chapter the architecture builds toward. The router of \Cref{ch:router}
sends every query it cannot place on a cheap tier to retrieval, as a fallback. That
fallback rests on an assumption that turns out to be false: that retrieval always
helps. The retrieval-utility classifier is \caem{}'s correction. It is a small
learned model that sits on the two fall-through paths and decides, per query,
whether retrieval will genuinely help the generator or quietly hurt it. It is the
newest component of the system and the one that most directly improves the
architecture's behaviour on the questions where ordinary retrieval-augmented
generation fails.

\begin{intuition}[Retrieval is a tool, not a reflex]
Retrieval-augmented generation is usually framed as strictly safer: ground the
model in evidence and it hallucinates less. That is true on the questions where the
evidence supports the right answer. It is false on a specific and important class of
questions, the ones built around a common misconception or the ones the model
already knows well, where the retrieved passages either reinforce the trap or
distract from a stronger parametric answer. On those, retrieval makes the answer
\emph{worse}. The retrieval-utility classifier learns to tell the two cases apart,
so that retrieval becomes a tool used when it helps rather than a reflex applied to
everything.
\end{intuition}

\section{The wrong-default failure mode}
\label{sec:ruc-motivation}

The empirical anchor for this chapter is a regression the architecture showed
\emph{before} the classifier existed. On misconception-bait benchmarks, exact-match
accuracy declined cycle over cycle. The cause was not the verifier and not the
memory; it was the unconditional retrieval dispatch feeding the self-improvement
loop. When a misconception question was sent to retrieval, the retrieved passages
nudged the generator toward the popular-but-wrong answer, the verifier sometimes let
that confident answer through, and the self-improvement loop of \Cref{ch:cycle} then
folded the retrieval-amplified error into the model's own weights. Each cycle made
the model a little more confidently wrong on exactly the questions designed to catch
it.

\begin{precise}[The tier-conditional self-improvement hypothesis]
The classifier's deeper purpose is to make self-improvement \emph{benchmark
conditional}. The self-improvement loop trains on the answers the system produced and
trusted, so the \emph{tier} that produced an answer determines what the model learns
from it. If retrieval-hurt queries are routed to retrieval, the loop ingests
retrieval-polluted answers and the model drifts on those benchmarks. If instead
those queries are routed to zero-shot generation, the loop ingests clean parametric
answers and the model improves. The classifier is therefore not only a
per-query cost optimiser; it is the gate that decides \emph{which kind of answer
enters the training pool} for each query. Routing retrieval-helpful queries to the
grounded tier and retrieval-hurt queries to the zero-shot tier lets the loop learn
grounded reasoning where grounding helps and clean parametric reasoning where it does
not, which is what allows one self-improvement loop to raise accuracy across a panel
of benchmarks that reward opposite things.
\end{precise}

\begin{bythenumbers}[The asymmetric-rescue payoff, with one caveat]
The deployed on/off ablation puts numbers on the mechanism. Turning the classifier on
lifts pooled exact match by about $5.8$ percentage points and collapses the
retrieval-tier share by about $34$ points, and the gains are concentrated exactly where
retrieval was hurting: the commonsense probe rises by about $14$ points and the
misconception probe (exact match judged by the language-model judge) by about $10$
points, while the retrieval-helpful fact-checking benchmark stays flat. The pattern is
the asymmetric rescue the hypothesis predicts: lift the benchmarks retrieval was hurting
without disturbing the one it was helping. One caveat keeps the contrast honest. This is
a bundled on/off comparison, with the classifier enabled alongside concurrent pipeline
corrections, not a clean classifier-only isolation, so the figures bound the joint
effect of the change rather than attribute every point to the classifier alone.
\end{bythenumbers}

\begin{figure}[htbp]
\centering
\begin{tikzpicture}[
  font=\small\sffamily,
  q/.style={rectangle, rounded corners=3pt, draw=cbInk!55, fill=cbInk!6,
    line width=0.9pt, align=center, inner sep=4pt, text width=3.0cm, font=\footnotesize},
  tA/.style={rectangle, rounded corners=4pt, draw=cbPitfall!75!black, fill=cbPitfall!10,
    line width=1pt, align=center, inner sep=4pt, text width=3.3cm, font=\footnotesize},
  tB/.style={rectangle, rounded corners=4pt, draw=cbExample!70!black, fill=cbExample!12,
    line width=1pt, align=center, inner sep=4pt, text width=3.3cm, font=\footnotesize},
  pool/.style={rectangle, rounded corners=4pt, draw=cbFormula!70!black, fill=cbFormula!10,
    line width=1pt, align=center, inner sep=4pt, text width=2.4cm, font=\footnotesize},
  flow/.style={-{Stealth[length=2.2mm]}, line width=1pt, draw=cbInk!70},
]
  % row 1: without RUC
  \node[font=\footnotesize\bfseries, text=cbPitfall!70!black, anchor=west] at (-0.2,3.3) {Without the classifier};
  \node[q]  (q1) at (0.6,2.3) {retrieval-hurt query\\(misconception bait)};
  \node[tA] (t1) at (4.6,2.3) {Tier 3 retrieval $\to$ passage-amplified \emph{wrong} answer};
  \node[pool] (p1) at (8.8,2.3) {SIL pool\\learns the error};
  \draw[flow] (q1) -- (t1);
  \draw[flow] (t1) -- (p1);
  \node[font=\scriptsize, text=cbPitfall!70!black, anchor=west] at (10.2,2.3) {model drifts};
  % row 2: with RUC
  \node[font=\footnotesize\bfseries, text=cbExample!55!black, anchor=west] at (-0.2,1.3) {With the classifier};
  \node[q]  (q2) at (0.6,0.3) {retrieval-hurt query\\(misconception bait)};
  \node[tB] (t2) at (4.6,0.3) {Tier 2 zero-shot $\to$ clean \emph{parametric} answer};
  \node[pool] (p2) at (8.8,0.3) {SIL pool\\learns clean reasoning};
  \draw[flow] (q2) -- (t2);
  \draw[flow] (t2) -- (p2);
  \node[font=\scriptsize, text=cbExample!55!black, anchor=west] at (10.2,0.3) {model improves};
\end{tikzpicture}
\caption[Tier-conditional self-improvement]{Why the classifier matters beyond cost.
On a retrieval-hurt query, the unconditional dispatch (top) sends a
passage-amplified wrong answer into the self-improvement pool, and the model drifts
on that benchmark cycle over cycle. The classifier (bottom) routes the same query to
zero-shot generation, so a clean parametric answer enters the pool and the model
improves. The classifier decides not just which tier answers, but which kind of
answer the loop learns from. (Schematic.)}
\label{fig:tier-conditional}
\end{figure}

\section{Where the classifier fires, and where it does not}
\label{sec:ruc-where}

The classifier is a refinement, not a replacement, so it is important to be exact
about its reach. It is consulted at exactly the two points in \Cref{ch:router} that
would otherwise route unconditionally to retrieval: the safety-override path and the
score fall-through. At both, the router's default is the grounded tier, and the
classifier is given the chance to redirect that default to zero-shot generation
instead.

\begin{pitfall}[The classifier never overrides a cheap-tier decision]
A query that earns Tier 1 on its combined score, or Tier 2 on its similarity, never
consults the classifier at all. The memory-direct and similarity-based zero-shot
decisions stand on their own. The classifier only ever chooses between the grounded
tier and the zero-shot tier \emph{for queries the router was about to send to
retrieval}. Its job is to split the retrieval fallback into ``really retrieve'' and
``actually, generate directly,'' and nothing more. This keeps its influence
localised and its failures bounded: a misfire can only move a query between the two
most expensive tiers, never onto the memory-direct path.
\end{pitfall}

Concretely, the classifier reads a feature vector for the query, produces a
probability $\prag$ that retrieval is the better path, and the router escalates to the
grounded tier if and only if that probability clears a threshold. When no classifier
is wired in, or the query text is unavailable, the router falls back to its old
behaviour and routes the fall-through to retrieval, so the classifier degrades safely
to the pre-classifier system.

\begin{figure}[htbp]
\centering
\begin{tikzpicture}[
  font=\small\sffamily,
  src/.style={rectangle, rounded corners=3pt, draw=cbInk!55, fill=cbInk!6,
    line width=0.9pt, align=center, inner sep=5pt, text width=3.2cm, font=\footnotesize},
  ruc/.style={rectangle, rounded corners=5pt, draw=cbFormula!75!black, fill=cbFormula!10,
    line width=1.2pt, align=center, inner sep=6pt, text width=3.0cm,
    drop shadow={shadow xshift=0.3ex, shadow yshift=-0.3ex, fill=black!10}},
  t3/.style={rectangle, rounded corners=4pt, draw=cbScenario!70!black, fill=cbScenario!12,
    line width=1pt, align=center, inner sep=5pt, text width=2.5cm, font=\footnotesize},
  t2/.style={rectangle, rounded corners=4pt, draw=cbIntuition!75!black, fill=cbIntuition!12,
    line width=1pt, align=center, inner sep=5pt, text width=2.5cm, font=\footnotesize},
  flow/.style={-{Stealth[length=2.4mm]}, line width=1.1pt, draw=cbInk!75},
]
  \node[src] (veto) at (0,1.0) {safety override\\($\upre < 0.38$)};
  \node[src] (fall) at (0,-1.0) {score fall-through\\(no cheap tier earned)};
  \node[ruc] (ruc) at (4.2,0) {\textbf{Retrieval-utility classifier}\\$\prag = \Pr(\text{RAG helps})$};
  \node[t3] (t3) at (9.6,1.0) {\textbf{Tier 3}\\retrieve + ground};
  \node[t2] (t2) at (9.6,-1.0) {\textbf{Tier 2}\\zero-shot};
  \draw[flow] (veto) -- (ruc);
  \draw[flow] (fall) -- (ruc);
  \draw[flow] (ruc) -- node[above=2pt, sloped, font=\scriptsize, text=cbScenario!60!black]{$\prag \ge 0.375$} (t3);
  \draw[flow] (ruc) -- node[below=2pt, sloped, font=\scriptsize, text=cbIntuition!65!black]{$\prag < 0.375$} (t2);
\end{tikzpicture}
\caption[Where the classifier sits]{The classifier intercepts the two router
outcomes that default to retrieval. It emits a probability that retrieval helps; the
router commits to the grounded tier above the threshold and to zero-shot generation
below it. Tier 1 and similarity-based Tier 2 decisions never reach this box. (Schematic.)}
\label{fig:ruc-where}
\end{figure}

\section{What the classifier is}
\label{sec:ruc-what}

For all its importance, the classifier is deliberately simple: a logistic regression
on a standardised feature vector. There is no neural network, no fine-tuning, and no
dependence on the generator's weights at inference. Thirty-six engineered features
are assembled for the query, standardised to zero mean and unit variance, and fed to
a logistic regression that outputs $\prag$. The decision is a threshold at
$\tauruc = 0.375$.

\begin{lstlisting}[caption={The classifier's prediction, condensed from \texttt{caem/routing/ruc.py}.}, label={lst:ruc-predict}]
feats = build_feature_vector(question,                 # 23 surface-form features
                             top1_passage_sim,         # + retrieval-evidence features
                             top1_passage_entity_overlap,
                             p_ik,                      # + self-knowledge probe
                             **extra_features)          # 36 features in total
x     = [feats[f] for f in feature_order]
x_s   = scaler.transform([x])                          # standardise (zero mean, unit var)
p_rag = lr.predict_proba(x_s)[0, 1]                    # P(retrieval is the better path)
decision = "RAG" if p_rag >= 0.375 else "DIRECT"       # tau_RUC = 0.375
\end{lstlisting}

\begin{intuition}[Why a logistic regression and not a neural network]
A more powerful model is tempting, but three reasons make the linear choice the right
one here. First, the decision is binary and the signal is concentrated in a handful
of features, so a linear boundary already separates the classes well. Second, a
logistic regression's coefficients are directly readable, which lets us audit
\emph{why} the classifier routes a query the way it does, a property we use in
\Cref{sec:ruc-pik}. Third, with only a few thousand labelled queries, a linear model
with strong regularisation generalises where a larger model would overfit. The
classifier is small on purpose; its leverage comes from the labels it was trained on,
not from its capacity.
\end{intuition}

\section{The features, by family}
\label{sec:ruc-features}

The thirty-six features fall into three families, distinguished by what they cost to
compute and what they capture. \Cref{tab:ruc-features} lays them out.

\begin{table}[htbp]
\centering
\caption[The classifier's feature families]{The thirty-six features, by family. The
surface-form family is computed inline from the question text; the retrieval-evidence
family is supplied by the pipeline from a cheap top-passage lookup; the self-knowledge
probe is read from the base model.}
\label{tab:ruc-features}
\small
\renewcommand{\arraystretch}{1.22}
\begin{tabularx}{\textwidth}{@{}l c >{\raggedright\arraybackslash}X@{}}
\toprule
\textbf{Family} & \textbf{Count} & \textbf{What it captures} \\
\midrule
\textbf{Question surface form}
  & 23
  & Cheap, text-only cues: question length, presence of negation, temporal,
    numerical, and misconception-bait phrasing; the interrogative type (who, what,
    when, where, why, how, yes-no); named-entity count and entity popularity;
    adversarial, opinion-seeking, and open-ended phrasing; and answer-type matches
    between the question and the top passage (does the question want a year or a
    number, and does the passage contain one). \\
\addlinespace[2pt]
\textbf{Retrieval evidence}
  & 12
  & Whether the corpus actually has good evidence: the top passage's similarity and
    entity overlap with the question; the top-five similarity spread (max, mean,
    standard deviation); cross-encoder rerank scores (top one, top-five mean and
    spread); pairwise entailment among the retrieved passages (mean, minimum,
    spread), which measures whether the passages agree; and a binary Wikidata
    entity hit. \\
\addlinespace[2pt]
\textbf{Self-knowledge}
  & 1
  & The probe $\pik$: the base model's own estimate of whether it can answer the
    question correctly without retrieval. This single feature carries the most
    weight of any in the model. \\
\bottomrule
\end{tabularx}
\end{table}

The three families answer three complementary questions. The surface-form family asks
\emph{what kind of question is this}, catching the misconception-bait and
opinion-seeking phrasings where retrieval tends to hurt. The retrieval-evidence family
asks \emph{does the corpus actually have something useful}, since retrieval cannot
help when the passages are irrelevant or disagree with one another. And the
self-knowledge probe asks \emph{does the model already know this}, because a question
the model can answer from its own parameters gains nothing from retrieval and may be
distracted by it.

\section{The star feature: the self-knowledge probe}
\label{sec:ruc-pik}

If one feature defines the classifier, it is the self-knowledge probe $\pik$. It is by
far the strongest, with a coefficient about two and a half times the next largest in
absolute value (about $-1.77$ on standardised features against about $+0.70$ for the
next feature), and its sign is exactly what the design intends.

\begin{figure}[htbp]
\centering
\begin{tikzpicture}
\begin{axis}[
  width=11.6cm, height=7.2cm,
  xbar, y axis line style={opacity=0},
  axis x line=bottom, x axis line style={cbInk!40},
  xlabel={logistic-regression coefficient (standardised features)},
  xmin=-2.0, xmax=1.0,
  enlarge y limits=0.06, bar width=8pt,
  ytick=data, font=\small\sffamily, tick label style={font=\scriptsize},
  symbolic y coords={numerical cue, interrog what, top5 sim std, interrog who,
    interrog when, rerank top5 mean, interrog yes/no, rerank top5 std,
    negation, self-knowledge},
  nodes near coords, nodes near coords style={font=\tiny, /pgf/number format/fixed,
    /pgf/number format/precision=2},
  xmajorgrids, grid style={cbInk!8},
]
  \addplot[draw=cbFormula!70!black, fill=cbFormula!25] coordinates {
    (-1.77,self-knowledge) (0.70,negation) (0.49,rerank top5 std)
    (-0.38,interrog yes/no) (0.37,rerank top5 mean) (-0.35,interrog when)
    (-0.28,interrog who) (0.24,top5 sim std) (-0.23,interrog what)
    (0.22,numerical cue)
  };
  \draw[cbInk!50, dashed] (axis cs:0,numerical cue) -- (axis cs:0,self-knowledge);
\end{axis}
\end{tikzpicture}
\caption[Top features by weight]{The ten features with the largest coefficients.
Negative coefficients (left) push toward zero-shot generation; positive coefficients
(right) push toward retrieval. The self-knowledge probe $\pik$ dominates, and its
strong negative weight encodes the intended rule: the more the model already knows,
the less it should retrieve. The negation cue carries the largest positive weight, with
the rerank-spread features close behind, escalating to retrieval when the corpus has
useful, varied evidence. (Learned coefficients from the deployed classifier.)}
\label{fig:ruc-coefs}
\end{figure}

\begin{precise}[What the probe is, and why the adapter is switched off]
The probe is a direct-correctness estimator in the style of
Kadavath et al.~\cite{kadavath2022language}: a small head reads the base model's
last-token hidden state on the question and predicts the probability that the model
would answer correctly. Its coefficient in the classifier is strongly negative
(about $-1.77$ on standardised features), which is exactly right: a high probe value
means the model already knows the answer, so retrieval is unlikely to help and the
query should go to zero-shot generation. On its own the probe reaches a mean
out-of-fold area under the ROC curve of about $0.77$, so it is usefully predictive but
imperfect, which is precisely why it is one feature among many rather than the whole
decision. One implementation detail is load-bearing:
the probe is computed with the low-rank adapter of \Cref{ch:cycle} \emph{disabled}, so
that it reads the frozen base model's self-knowledge on the same distribution the
probe was trained on, rather than the shifting fine-tuned distribution. Measuring the
base model's knowledge, not the adapter's, keeps the signal stable across cycles.
\end{precise}

\section{How the classifier was trained}
\label{sec:ruc-training}

The classifier's leverage comes from its labels, so the labelling deserves care. The
question for every training query is not whether the answer is right, but whether
\emph{retrieval} produced a better answer than \emph{zero-shot generation}. To label a
query, both paths are run: a zero-shot answer and a retrieval-augmented answer are each
scored through the verifier of \Cref{ch:verifier}. The query is labelled positive
(retrieval helps) when the retrieval answer wins and negative (retrieval hurts) when the
zero-shot answer wins.

\begin{pitfall}[Train only on the disagreements]
Most queries are ties: both paths get it right, or both get it wrong. Those carry no
signal about \emph{which path to choose}, and including them would swamp the classifier
with uninformative examples and waste expensive verifier calls. So the training set is
filtered to the queries where the two paths \emph{disagree}, one right and one wrong.
This is both a statistical choice, concentrating the model on the decision boundary that
matters, and a cost choice, since only the disagreements need the double verifier
scoring. The shipped classifier was trained on roughly $3{,}200$ such disagreement
queries.
\end{pitfall}

\begin{bythenumbers}[Training and validation]
\begin{itemize}[leftmargin=1.3em, itemsep=1pt]
  \item Training set: $3{,}208$ disagreement queries ($1{,}846$ retrieval-helps,
  $1{,}362$ retrieval-hurts) drawn from \emph{eight} benchmarks, deliberately broader
  than the five-benchmark \caem{} panel so the classifier learns a transferable
  retrieval-utility signal.
  \item Model selection: nested cross-validation, five outer folds and five inner
  folds, with the regularisation strength and the decision threshold both chosen on
  held-out inner folds (final strength $C = 0.5$, threshold $\tauruc = 0.375$).
  \item Discrimination: pooled out-of-fold area under the ROC curve of $0.850$, with
  per-benchmark values ranging from $0.69$ on OpenBookQA to $0.94$ on TriviaQA.
\end{itemize}
\end{bythenumbers}

The nested cross-validation matters because two things are being chosen at once, the
regularisation strength and the threshold, and choosing either on the data used to
report performance would inflate the result. The outer folds give an honest estimate;
the inner folds do the choosing. The threshold landed below one half, at $0.375$,
because the cost of the two errors is asymmetric and the sweep that maximised net
utility preferred to retrieve a little more readily than a symmetric cutoff would.

\section{Honest limitations}
\label{sec:ruc-threats}

Two caveats keep the classifier in proportion.

\begin{pitfall}[Not every coefficient sign matches the naive intuition]
A registered audit checks whether the learned coefficients carry the signs the design
expects. The self-knowledge probe passes cleanly, and the rerank score passes, but the
top-passage similarity does not: its coefficient is mildly negative where a naive reading
would expect positive, so the audit records an overall fail on that check. The likely
reason is correlation among the retrieval features, with the richer rerank and
passage-agreement signals absorbing the useful part of raw similarity and leaving its
residual coefficient small and oppositely signed. It is reported transparently rather
than hidden, and it does not change the deployed behaviour, since discrimination is
measured on held-out folds, not on coefficient signs. It is a reminder that a linear
model's individual weights are interpretable only up to the correlations among its
inputs.
\end{pitfall}

\begin{pitfall}[Discrimination is uneven across benchmarks]
The pooled discrimination of $0.850$ hides real variation. The classifier separates the
two classes very well on TriviaQA and HaluEvalQA, and only modestly on FEVER, OpenBookQA,
and CommonsenseQA, where the retrieval-helps and retrieval-hurts queries look more alike
on these features. The classifier is therefore most valuable on the open-domain factoid
and misconception-bait benchmarks and least decisive on the short, self-contained ones,
which is consistent with where the wrong-default failure mode bites hardest.
\end{pitfall}

\section{The classifier in the lineage of selective retrieval}
\label{sec:ruc-lineage}

\begin{defence}[If an examiner asks: ``how is this different from prior selective-retrieval work?'']
The idea that retrieval should be conditional is not new, and the thesis does not claim
it. A line of work reaches the same diagnosis: that an unconditional retrieval dispatch
leaves accuracy on the table relative to a per-query routing decision, including
when-not-to-retrieve on PopQA~\cite{mallen2023whennot}, self-knowledge
routing~\cite{wang2023skr}, Adaptive-RAG~\cite{jeong2024adaptive}, and
Probing-RAG~\cite{baek2025probing}. \caem{}'s classifier sits in that lineage with two
distinguishing features. It is positioned \emph{downstream} of a memory-coverage and a
readiness mechanism, so it refines only the residual queries those mechanisms could not
place, rather than gating every query. And its labels are tied to \caem{}'s own
verifier-scored disagreements between the two paths, so it learns retrieval utility as
\emph{this} system experiences it, and feeds directly into the tier-conditional
self-improvement of \Cref{sec:ruc-motivation}, which the selective-retrieval line does
not address because those systems do not have a self-improvement loop to protect.
\end{defence}

\section{The two readings, refined}
\label{sec:ruc-scenarios}

Both of our running queries pass through the classifier, and it sends them opposite ways.

\begin{scenario}[The FEVER claim: redirected to zero-shot]
The claim ``The French Revolution began before 1792'' reaches the classifier through the
score fall-through on the empty cycle-zero memory. The query is short, self-contained, and
the kind of fact the base model knows, so the self-knowledge probe reads moderately high
and the surface-form features mark it as a simple declarative claim. The classifier returns
a low $\prag$, below $0.375$, and the router redirects the default retrieval dispatch to
zero-shot generation. The model labels the claim from its own knowledge, which is the right
call and avoids dragging in passages that could only muddy a fact it already holds.
\end{scenario}

\begin{scenario}[The TriviaQA question: confirmed for retrieval]
The question ``\,`If I Ruled The World' comes from which musical?'' reaches the classifier
through the safety override. Here the self-knowledge probe reads low, the model does not
confidently know the answer, and the retrieval-evidence features show that the corpus holds
relevant, agreeing passages. The classifier returns a high $\prag$, above $0.375$, and
confirms the grounded dispatch. Retrieval is genuinely the better path here, and the
classifier agrees with the safety override rather than countermanding it.
\end{scenario}

These two cases are the classifier's whole job in miniature: send the question the model
already knows to its own parameters, and send the question it does not know to the
evidence. \Cref{ch:tiers} now takes up Stage 4, where the chosen tier finally runs and
produces an answer to be judged.
